%=============================================================================
% Wave 3, Iteration 5: Cross-Reference Update
% Patent 14 (ICC / Cosmological Coupling) + Patent 15 (Generator / SRF)
%
% Agent 14/15 -- T1 Resolution via EM-Only Sourcing,
%                P15 M_eff-Corrected Detection Reassessment
% Date: March 2026
%
% W3i5 MANDATE:
%   P14: Attempt to RESOLVE T1 by the P14 critical caveat route:
%        if psi is sourced ONLY by EM stress-energy (as the DFD field
%        equation actually states), compute psi_cosmo from EM sources only.
%   P15: With M_eff = 3.01e6 kg correction, reassess all detection configs.
%
% Building on: W3i4_P14_update.tex (T1 genuine falsification, T2 inconsistency)
%              W3i4_P15_update.tex (material FOM table, Nb3Sn+YBCO confirmed)
%              W3i4_P11_update.tex (M_eff = 3.01e6 kg resolution)
%=============================================================================

\documentclass[11pt]{article}
\usepackage{amsmath,amssymb,amsthm}
\usepackage{geometry}
\geometry{margin=1in}
\usepackage{booktabs}
\usepackage{enumitem}
\usepackage{hyperref}
\usepackage{xcolor}
\usepackage{tcolorbox}
\usepackage{multirow}
\usepackage{array}
\usepackage{float}

\newtheorem{theorem}{Theorem}
\newtheorem{proposition}{Proposition}
\newtheorem{lemma}{Lemma}
\newtheorem{finding}{Finding}
\newtheorem{correction}{Correction}
\newtheorem{definition}{Definition}
\newtheorem{corollary}{Corollary}

\newcommand{\ee}[1]{\times 10^{#1}}
\newcommand{\psifield}{\psi}
\newcommand{\psib}{\bar{\psi}}
\newcommand{\GN}{G_N}
\newcommand{\Geff}{G_{\rm eff}}
\newcommand{\Gdot}{\dot{G}/G}
\newcommand{\psicosmo}{\bar{\psi}_{\rm cosmo}}
\newcommand{\aM}{a_0}
\newcommand{\DFD}{\textsc{dfd}}
\newcommand{\GR}{\textsc{gr}}
\newcommand{\LLR}{\textsc{llr}}
\newcommand{\EHT}{\textsc{eht}}
\newcommand{\RH}{R_H}
\newcommand{\Hzero}{H_0}
\newcommand{\Mpsi}{M_\psi}
\newcommand{\Meff}{M_\psi^{\rm eff}}
\newcommand{\lam}{\lambda}
\newcommand{\lamdiff}{|\lambda-1|}
\newcommand{\Opsi}{\Omega_\psi}
\newcommand{\gpsi}{\gamma_\psi}
\newcommand{\Vcav}{V_{\rm cav}}
\newcommand{\muG}{\mu_0^{\rm gal}}
\newcommand{\FOM}{\mathrm{FOM}}
\newcommand{\lambdaL}{\lambda_{\rm L}}
\newcommand{\xsc}{\xi_{\rm sc}}
\newcommand{\Tc}{T_{\rm c}}
\newcommand{\kB}{k_{\rm B}}
\newcommand{\order}[1]{\mathcal{O}(#1)}

%=============================================================================
\title{Wave 3 Iteration 5: Patent 14 (ICC) + Patent 15 (Generator)\\
\large T1 Resolution via EM-Only $\psi$-Sourcing:\\
Cosmological $\psi$-Background from Electromagnetic Sources,\\
$\dot{G}/G$ Reassessment, MOND without Mach's Principle,\\
and P15 $M_{\rm eff}$-Corrected Detection Roadmap}
\author{DFD Research Programme --- Agent 14/15 (W3i5)}
\date{March 2026}
%=============================================================================

\begin{document}
\maketitle
\tableofcontents
\newpage

%=============================================================================
\section*{W3i5 Context: The Most Critical Iteration}
%=============================================================================

\noindent\textbf{Status inherited from W3i4.}
W3i4 declared T1 ($\Gdot = +2.6\ee{-11}$~yr$^{-1}$, exceeding LLR by $347\times$)
a \emph{genuine falsification} of the minimal DFD cosmological coupling. However,
W3i4 also identified a critical caveat (T2): the DFD psi-field equation
\begin{equation}
\Box\psi + (1+\kappa)\left[(\nabla\psi)^2 - \dot{\psi}^2/c^2\right]
= \frac{4\pi G}{c^4} T_{\rm EM}
\label{eq:field-eq}
\end{equation}
is sourced \emph{only} by electromagnetic stress-energy $T_{\rm EM}$, not by
matter density $\rho_{\rm matter}$. The Mach integral that generates
$\psicosmo \approx 0.047$ assumed $\psi = GM/c^2r$ (the Newtonian
potential from \emph{matter}), which is inconsistent with Eq.~(\ref{eq:field-eq}).

\textbf{W3i5 mandate}: Follow the T2 caveat to its logical conclusion.
If $\psi$ is sourced only by EM stress-energy, what is the actual cosmological
$\psi$-background from electromagnetic sources? Does T1 resolve? What happens
to MOND, the shadow prediction, the Hubble tension, and S$_8$?

For P15: with the W3i4 P11 correction $\Meff = 3.01\ee{6}$~kg replacing
the placeholder $\Mpsi = 10^{-6}$~kg, reassess all P15 detection
configurations and determine which reaches survive.

%=============================================================================
\section{P14: The EM-Only $\psi$-Background}
\label{sec:EM-only}
%=============================================================================

\subsection{Electromagnetic energy density in the universe}

If $\psi$ is sourced exclusively by the electromagnetic stress-energy tensor
$T_{\rm EM}^{\mu\nu}$, then the cosmological $\psi$-background must be
computed from the EM energy density of the universe, not from matter.
We enumerate all significant cosmological EM sources.

\subsubsection{CMB radiation field}

The CMB has energy density:
\begin{equation}
u_{\rm CMB} = a_{\rm rad}\,T_{\rm CMB}^4
= \frac{\pi^2 k_B^4}{15\hbar^3 c^3}\,T_{\rm CMB}^4
= 4.17\ee{-14}~\text{J/m}^3
\label{eq:u-CMB}
\end{equation}
at $T_{\rm CMB} = 2.725$~K. The corresponding mass-equivalent density:
\begin{equation}
\rho_{\rm CMB} = \frac{u_{\rm CMB}}{c^2} = \frac{4.17\ee{-14}}{9\ee{16}}
= 4.63\ee{-31}~\text{kg/m}^3.
\label{eq:rho-CMB}
\end{equation}

However, for a pure radiation field, $T^{\rm EM}_{\mu}{}^{\mu} = 0$ (the
electromagnetic stress-energy tensor is traceless). This is a fundamental
property of the Maxwell field in 4 dimensions: $T^{\rm EM}_{\mu\nu}$
satisfies $g^{\mu\nu}T^{\rm EM}_{\mu\nu} = 0$.

\begin{finding}[CMB Does Not Source $\psi$ via the Trace]
\label{find:CMB-traceless}
The CMB photon field has a traceless stress-energy tensor:
$T^{\rm CMB}_{\mu}{}^{\mu} = 0$. Since the DFD field equation
(\ref{eq:field-eq}) couples $\psi$ to $T_{\rm EM}$ (which, for
Process~A sourcing, enters through the trace of the matter
stress-energy induced by EM fields in material), a free-space
CMB radiation field produces \emph{zero} direct $\psi$-source.

The CMB sources $\psi$ only indirectly, through its interaction with
matter (inverse Compton scattering, free-free absorption in the IGM),
which creates a non-zero trace through the matter kinetic energy response.
\end{finding}

\subsubsection{Indirect CMB sourcing via IGM Thomson scattering}

The CMB interacts with free electrons in the intergalactic medium (IGM)
via Thomson scattering. The scattered radiation transfers momentum to
electrons, creating a kinetic energy density in the IGM. The power
transferred per electron:
\begin{equation}
P_{\rm Thomson} = \sigma_T c\,u_{\rm CMB} = 6.65\ee{-29}\times3\ee{8}
\times 4.17\ee{-14} = 8.3\ee{-34}~\text{W}.
\end{equation}
The IGM electron density at $z=0$ is $n_e \approx 2\ee{-7}$~cm$^{-3}
= 0.2$~m$^{-3}$. The kinetic energy density transferred over the
Hubble time $t_H \approx 4.4\ee{17}$~s:
\begin{equation}
u_{\rm kin}^{\rm IGM} = n_e\,P_{\rm Thomson}\,t_H
= 0.2\times8.3\ee{-34}\times4.4\ee{17}
= 7.3\ee{-17}~\text{J/m}^3.
\label{eq:u-kin-IGM}
\end{equation}
This is a factor $\sim 600$ below $u_{\rm CMB}$ itself. The effective
stress-energy trace from this kinetic contribution is:
\begin{equation}
T^{\rm IGM}_{\mu}{}^{\mu} \sim u_{\rm kin}^{\rm IGM}
\sim 7.3\ee{-17}~\text{J/m}^3.
\end{equation}

\subsubsection{Galactic magnetic fields}

The Milky Way has a coherent magnetic field $B_{\rm MW} \approx 3$--$6~\mu$G
$= 3$--$6\ee{-10}$~T. The EM energy density:
\begin{equation}
u_{\rm MW} = \frac{B_{\rm MW}^2}{2\mu_0}
= \frac{(5\ee{-10})^2}{2\times1.26\ee{-6}}
= 9.9\ee{-14}~\text{J/m}^3.
\label{eq:u-MW}
\end{equation}

Again, the free-field $B$-field has a traceless stress-energy tensor.
The $\psi$-source arises from the interaction of $B_{\rm MW}$ with
the ISM plasma. The magnetised plasma has a non-zero stress-energy
trace through magnetic pressure acting on the gas:
\begin{equation}
T^{\rm ISM}_{\mu}{}^{\mu} = \rho_{\rm ISM}c^2 - 3P_{\rm ISM}
\approx \rho_{\rm ISM}c^2\left(1 - \frac{3P_{\rm ISM}}{\rho_{\rm ISM}c^2}\right).
\end{equation}
For non-relativistic ISM gas, $P/(\rho c^2) \sim v_{\rm th}^2/c^2
\sim 10^{-10}$, so $T_{\mu}{}^{\mu} \approx \rho_{\rm ISM}c^2$.
But this is the \emph{matter} trace, not the EM trace. The pure EM
contribution to the trace remains zero.

\begin{finding}[Pure EM Fields Do Not Source $\psi$ via the Trace Coupling]
\label{find:EM-traceless}
All free electromagnetic fields (radiation, static $B$-fields, $E$-fields)
have $T^{\rm EM}_{\mu}{}^{\mu} = 0$ in 4 dimensions. The DFD field equation
sourced by $T_{\rm EM}$ cannot generate a cosmological $\psi$-background
from free EM fields. The only EM-to-$\psi$ sourcing occurs through
Process~A: EM fields interacting with \emph{matter} (cavity walls,
plasma, ISM gas), which creates a material stress-energy trace.
\end{finding}

\subsubsection{Intergalactic magnetic fields}

Intergalactic $B$-fields are constrained to $B_{\rm IGM} \lesssim 10^{-9}$~G
$= 10^{-13}$~T by Faraday rotation measures and CMB constraints. The
energy density:
\begin{equation}
u_{\rm IGM}^B = \frac{(10^{-13})^2}{2\times1.26\ee{-6}}
= 4.0\ee{-21}~\text{J/m}^3.
\label{eq:u-IGM-B}
\end{equation}
This is 7 orders below even the indirect IGM Thomson contribution.
The $\psi$-source from IGM $B$-fields interacting with the diffuse
IGM plasma is negligible.

\subsection{The EM-only $\psicosmo$: quantitative estimate}

The dominant EM-sourced $\psi$-contribution comes from the indirect
channel: CMB photons scattering off IGM electrons, producing a kinetic
energy density in the IGM that has a non-zero stress-energy trace.

Using the Mach-integral approach but now with $\rho \to \rho_{\rm EM-eff}
= u_{\rm kin}^{\rm IGM}/c^2$:
\begin{equation}
\psicosmo^{\rm EM} = \frac{4\pi G\,\rho_{\rm EM-eff}\,R_H^2}{c^2}
= \frac{4\pi\times6.67\ee{-11}\times\frac{7.3\ee{-17}}{9\ee{16}}
\times(1.32\ee{26})^2}{9\ee{16}}.
\label{eq:psi-EM-calc}
\end{equation}
Computing:
\begin{align}
\rho_{\rm EM-eff} &= 8.1\ee{-34}~\text{kg/m}^3, \\
4\pi G\,\rho_{\rm EM-eff} &= 4\pi\times6.67\ee{-11}\times8.1\ee{-34}
= 6.8\ee{-43}~\text{s}^{-2}, \\
R_H^2 &= 1.74\ee{52}~\text{m}^2, \\
\psicosmo^{\rm EM} &= \frac{6.8\ee{-43}\times1.74\ee{52}}{9\ee{16}}
= \frac{1.18\ee{10}}{9\ee{16}}
= 1.3\ee{-7}.
\label{eq:psi-EM-value}
\end{align}

\begin{tcolorbox}[colback=green!5!white, colframe=green!60!black,
  title=\textbf{KEY RESULT: $\psicosmo^{\rm EM} \approx 1.3\ee{-7}$}]
If the DFD $\psi$-field is sourced only by electromagnetic stress-energy
(as the field equation states), the cosmological $\psi$-background from
all EM sources is:
\begin{equation}
\boxed{\psicosmo^{\rm EM} \approx 1.3\ee{-7}}
\end{equation}
This is a factor of $3.6\ee{5}$ smaller than the Mach-integral value
$\psicosmo^{\rm matter} \approx 0.047$. The dominant contribution is
indirect: CMB Thomson scattering in the IGM.

\textbf{Upper bound}: Including all galactic magnetic field contributions
(integrated over $\sim 10^{11}$ galaxies, each with $u_{\rm gal} \sim
10^{-13}$~J/m$^3$ in a volume $\sim (30~\text{kpc})^3$), the total
galactic EM contribution to $\psicosmo$ adds $\lesssim 10^{-8}$.
The conservative upper bound is:
\begin{equation}
\psicosmo^{\rm EM} < 2\ee{-7}.
\end{equation}
\end{tcolorbox}

%=============================================================================
\section{T1 Resolution: $\Gdot$ from EM-Only Sourcing}
\label{sec:T1-resolution}
%=============================================================================

\subsection{Revised $\Gdot$ prediction}

With $\psicosmo^{\rm EM} \approx 1.3\ee{-7}$ replacing $\psicosmo^{\rm matter}
\approx 0.047$, the $\Gdot$ prediction from the Mach-integral route becomes:
\begin{align}
\frac{\dot{G}}{G}\bigg|_{\rm EM-only}
&= 2\psicosmo^{\rm EM}\,H_0(5+2q_0) \nonumber\\
&= 2\times1.3\ee{-7}\times2.27\ee{-18}\times(5+2\times(-0.55))
\nonumber\\
&= 2\times1.3\ee{-7}\times2.27\ee{-18}\times3.9 \nonumber\\
&= 2.3\ee{-24}~\text{s}^{-1} \nonumber\\
&= 7.3\ee{-17}~\text{yr}^{-1}.
\label{eq:Gdot-EM}
\end{align}

\begin{tcolorbox}[colback=green!5!white, colframe=green!60!black,
  title=\textbf{T1 RESOLVED: $\Gdot$ from EM-Only Sourcing Is 1000$\times$ Below LLR}]
\begin{equation}
\boxed{\frac{\dot{G}}{G}\bigg|_{\rm EM-only} \approx 7.3\ee{-17}~\text{yr}^{-1}}
\end{equation}
This is a factor of $\sim 1000$ \emph{below} the LLR bound
($< 7.5\ee{-14}$~yr$^{-1}$) and a factor of $\sim 22$ below the
planetary ephemeris bound ($< 3\ee{-13}$~yr$^{-1}$).

\textbf{The $\dot{G}/G$ tension T1 is completely resolved if the DFD
psi-field equation is taken at face value}: $\psi$ is sourced only by
EM stress-energy, not by matter. The cosmological $\psi$-background
from EM sources alone is too small to produce a detectable $\dot{G}/G$.
\end{tcolorbox}

\subsection{Comparison with all observational bounds}

\begin{table}[h]
\centering
\caption{$\Gdot$ bounds vs.\ the EM-only DFD prediction.}
\label{tab:Gdot-EM}
\begin{tabular}{llll}
\toprule
\textbf{Method} & \textbf{Bound (yr$^{-1}$)} & \textbf{DFD EM-only} & \textbf{Status} \\
\midrule
Lunar laser ranging & $< 7.5\ee{-14}$ & $7.3\ee{-17}$ & \textbf{1000$\times$ below} \\
Planetary ephemeris & $< 3\ee{-13}$ & $7.3\ee{-17}$ & \textbf{4100$\times$ below} \\
Helioseismology & $< 1.6\ee{-12}$ & $7.3\ee{-17}$ & \textbf{22000$\times$ below} \\
Millisecond pulsars & $< 1.1\ee{-12}$ & $7.3\ee{-17}$ & \textbf{15000$\times$ below} \\
Binary pulsar & $< 2.0\ee{-12}$ & $7.3\ee{-17}$ & \textbf{27000$\times$ below} \\
White dwarf cooling & $< 1.8\ee{-11}$ & $7.3\ee{-17}$ & \textbf{$2.5\ee{5}\times$ below} \\
BBN & $< 4\ee{-12}$ & $7.3\ee{-17}$ & \textbf{55000$\times$ below} \\
\midrule
\textbf{DFD (EM-only)} & --- & $\mathbf{7.3\ee{-17}}$ & \textbf{All satisfied} \\
\textbf{DFD (Mach, W3i4)} & --- & $2.6\ee{-11}$ & Excluded $347\times$ \\
\bottomrule
\end{tabular}
\end{table}

The EM-only interpretation satisfies \emph{every} $\Gdot$ bound by orders
of magnitude. This is a qualitative change from the W3i4 finding of
``genuine falsification.''

\subsection{The topological route: also resolved}

W3i4 showed the topological invariant $G\hbar H_0^2/c^5 = \alpha^{57}$
gives $\Gdot = +6.4\ee{-11}$~yr$^{-1}$ if $G$ tracks $H^{-2}$ at all
epochs. However, with the EM-only interpretation, $G_{\rm eff}$ receives
only the tiny correction $e^{-2\psicosmo^{\rm EM}} \approx 1 - 2.6\ee{-7}$.
The topological relation then holds as a \emph{static} (attractor)
relation at the current epoch, with $\dot{G}/G$ dominated by the EM-only
contribution. The $853\times$ violation from W3i4 disappears.

%=============================================================================
\section{Consequences: What Is Lost}
\label{sec:lost}
%=============================================================================

\subsection{Mach's principle: lost in its strong form}

The DFD Mach integral $\psicosmo = 4\pi G\bar{\rho}_{\rm matter}R_H^2/c^2
\approx 0.047$ was the foundation of the DFD claim to implement Mach's
principle: inertia arises from the gravitational psi-background of all
matter in the universe, with $m_{\rm DFD} = m_0\,e^{\psicosmo}$.

If $\psi$ is sourced only by EM and $\psicosmo^{\rm EM} \approx 1.3\ee{-7}$,
then:
\begin{equation}
m_{\rm DFD} = m_0\,e^{\psicosmo^{\rm EM}} \approx m_0(1 + 1.3\ee{-7}).
\end{equation}
The ``inertial correction from the universe'' is $\sim 10^{-7}$---negligible.
Mach's principle in its strong form (inertia arises from distant matter)
is \emph{not implemented} by the EM-only DFD.

\begin{finding}[Mach's Principle Abandoned in EM-Only DFD]
\label{find:mach-lost}
The EM-only interpretation resolves T1 at the cost of abandoning the strong
Machian programme. Inertia does \emph{not} arise from the cosmological
psi-background in the EM-only version. The inertial mass correction from
$\psicosmo^{\rm EM}$ is $\sim 10^{-7}$, not the $\sim 5\%$ correction
from the Mach-integral value. This is an honest trade-off: T1 resolution
vs.\ Mach's principle.
\end{finding}

\subsection{The DFD-MOND connection: requires reformulation}

\subsubsection{The P5 MOND formula}

P5 derived the DFD MOND acceleration scale as:
\begin{equation}
a_0^{\rm DFD} = c\,H_0\,\psicosmo.
\label{eq:a0-DFD}
\end{equation}
With $\psicosmo^{\rm EM} \approx 1.3\ee{-7}$:
\begin{equation}
a_0^{\rm EM} = 3\ee{8}\times2.27\ee{-18}\times1.3\ee{-7}
= 8.9\ee{-17}~\text{m/s}^2.
\label{eq:a0-EM}
\end{equation}
This is a factor of $1.3\ee{6}$ below the observed $a_0 = 1.2\ee{-10}$~m/s$^2$.

\begin{finding}[MOND Scale Cannot Come from $\psicosmo^{\rm EM}$]
\label{find:MOND-lost}
The observed MOND acceleration scale $a_0 = 1.2\ee{-10}$~m/s$^2$ is
$1.3\ee{6}$ times larger than what the EM-only $\psicosmo$ can produce
via Eq.~(\ref{eq:a0-DFD}). The P5 formula $a_0 = cH_0\psicosmo$ is
\emph{not viable} in the EM-only framework.
\end{finding}

\subsubsection{Can $\mu(x)$ emerge without Mach's principle?}

The MOND interpolating function $\mu(x) = x/\sqrt{1+x^2}$ (with
$x = a/a_0$) in DFD emerges from the nonlinear self-coupling term
$(1+\kappa)(\nabla\psi)^2$ in the field equation (\ref{eq:field-eq}).
In the \emph{Newtonian} limit where $\psi = GM/c^2r$, this nonlinearity
creates a transition between Newtonian ($a \gg a_0$) and deep-MOND
($a \ll a_0$) regimes.

The nonlinear self-coupling is a property of the field equation
\emph{independent} of the source. It operates on any $\psi$-gradient,
regardless of whether the gradient was sourced by matter (Mach) or EM.
However, the \emph{scale} at which the transition occurs depends on
$\psicosmo$:
\begin{equation}
x = \frac{|\nabla\psi|\,c^2}{2\psicosmo\,H_0\,c}
= \frac{a}{a_0^{\rm DFD}},
\end{equation}
so $a_0$ is set by the cosmological background.

\textbf{Resolution path}: The MOND scale in DFD need not come from
$\psicosmo$ if the nonlinear self-coupling has a fixed point or an
intrinsic scale. The $(1+\kappa)(\nabla\psi)^2$ term has a natural
transition at:
\begin{equation}
|\nabla\psi|_{\rm crit} = \frac{1}{1+\kappa}
\label{eq:psi-crit}
\end{equation}
where the nonlinear term becomes comparable to the linear d'Alembertian.
If $\kappa$ is determined by the DFD topological invariant or the
coupling constant $\lambda$, this could set $a_0$ independently of
$\psicosmo$.

\begin{finding}[Alternative MOND Mechanism: Self-Coupling Fixed Point]
\label{find:MOND-alt}
The DFD nonlinear self-coupling $(1+\kappa)(\nabla\psi)^2$ has an
intrinsic transition scale $|\nabla\psi|_{\rm crit} = (1+\kappa)^{-1}$.
The corresponding acceleration scale is:
\begin{equation}
a_0^{\rm self} = \frac{c^2}{2(1+\kappa)}.
\end{equation}
To match $a_0 = 1.2\ee{-10}$~m/s$^2$:
\begin{equation}
(1+\kappa) = \frac{c^2}{2a_0} = \frac{9\ee{16}}{2.4\ee{-10}}
= 3.75\ee{26}.
\end{equation}
This requires $\kappa \approx 3.75\ee{26}$---an enormously large
nonlinear coupling constant. This is physically unappealing but
mathematically consistent. Whether this can be derived from the DFD
topological structure remains an open question for subsequent iterations.
\end{finding}

\subsection{P14+P5 incompatibility: dissolved by construction}

The W3i4 P14+P5 incompatibility ($\psicosmo = 0.047$ vs.\ $0.176$,
a $3.7\times$ mismatch) arose because the Mach integral and the MOND
formula demanded different values of $\psicosmo$. In the EM-only
framework, neither the Mach integral nor the P5 MOND formula is valid
in its original form. The incompatibility is dissolved---but replaced
by the open problem of deriving $a_0$ from a different mechanism.

%=============================================================================
\section{Consequences: What Survives and Is Strengthened}
\label{sec:survives}
%=============================================================================

\subsection{Shadow $+4.6\%$: survives and is strengthened}

The shadow prediction $\theta_{\rm DFD} = 1.046\,\theta_{\rm GR}$
arises from coherent L/T mode interference in the psi-screen around
the black hole. This effect depends on $\psi(r)$ near the photon sphere
($\psi = 1/3$), which is sourced by the local BH mass. It does
\emph{not} depend on $\psicosmo$.

In fact, the EM-only interpretation \emph{strengthens} the shadow
prediction: W3i4 noted a $\sim 5\%$ systematic in mass calibration
from the $G_{\rm eff}$ correction $e^{-2\times 0.047} \approx 0.91$.
With $\psicosmo^{\rm EM} \approx 1.3\ee{-7}$, the correction is
$e^{-2\times 1.3\ee{-7}} \approx 1 - 2.6\ee{-7}$, completely negligible.
The mass calibration systematic disappears, and the $+4.6\%$ prediction
is clean.

\begin{finding}[Shadow $+4.6\%$ Strengthened by EM-Only Interpretation]
\label{find:shadow-strengthened}
The shadow $+4.6\%$ prediction is \emph{independent} of $\psicosmo$ and
survives the EM-only interpretation. The W3i4 mass-calibration systematic
($\sim 5\%$ from $G_{\rm eff}$ correction) is eliminated: with
$\psicosmo^{\rm EM} \approx 10^{-7}$, the $G_{\rm eff}$ correction is
$< 10^{-6}$, far below any measurement precision. The $+4.6\%$ prediction
is cleaner in the EM-only framework.
\end{finding}

\subsection{Hubble tension: partially lost}

The DFD Hubble tension explanation (W3i2: $\Delta H_0 = 3.9 \pm 1.1$~km/s/Mpc,
70\% of the tension) derives from MOND-enhanced void outflow. The MOND
enhancement depends on $\mu(x)$ transitioning to the deep-MOND regime
at $a \lesssim a_0$ in low-density voids.

If the MOND scale $a_0$ is preserved by the self-coupling mechanism
(Finding~\ref{find:MOND-alt}), the Hubble tension explanation survives
in full. If $a_0$ cannot be derived independently, the Hubble tension
explanation is lost.

\textbf{Status}: Conditional on the self-coupling MOND mechanism.

\subsection{S$_8$ tension prediction: conditional}

The DFD S$_8$ prediction ($S_8^{\rm DFD} \sim 0.77$--$0.79$) depends on
the scale-dependent $G_{\rm eff}(k)$ from $\mu(x_k)$. This again depends
on the MOND interpolating function, which requires $a_0$.

\textbf{Status}: Same conditionality as the Hubble tension.

\subsection{DESI DR2 tension: relaxed}

The W3i2 DESI DR2 tension (2.4 sigma in $w_0$) arose from the psi-screen
refractive bias mechanism with $\psicosmo \approx 0.047$. With
$\psicosmo^{\rm EM} \approx 1.3\ee{-7}$, the refractive bias is
reduced by a factor $\sim 3.6\ee{5}$, making it negligible. The DFD
dark energy prediction via psi-screen bias effectively \emph{vanishes}
in the EM-only framework, and there is no prediction to be in tension
with DESI DR2.

However, if the dark energy replacement mechanism is refractive bias,
then $\psicosmo^{\rm EM} \approx 10^{-7}$ is far too small to produce
the observed apparent acceleration. The dark energy explanation via
accumulated refractive bias (Rank~3 in the Master Findings) is
\emph{lost} in the EM-only framework.

\begin{tcolorbox}[colback=yellow!5!white, colframe=orange!60!black,
  title=\textbf{EM-Only Trade-Off Summary}]
\textbf{Resolved:}
\begin{itemize}[nosep]
\item T1 ($\Gdot$): $347\times$ violation $\to$ $1000\times$ \emph{below} LLR
\item T2 (internal inconsistency): dissolved---field equation self-consistent
\item P14+P5 incompatibility: dissolved---both formulas invalid in old form
\item DESI DR2 tension: dissolved (no prediction to be in tension)
\item Shadow mass-calibration systematic: eliminated
\end{itemize}

\textbf{Lost:}
\begin{itemize}[nosep]
\item Mach's principle (strong form): inertial correction $\sim 10^{-7}$
\item MOND scale from cosmological background: requires alternative mechanism
\item Dark energy as psi-screen bias: $\psicosmo$ too small by $10^5$
\item Hubble tension explanation: conditional on alternative MOND
\item S$_8$ prediction: conditional on alternative MOND
\end{itemize}

\textbf{Strengthened:}
\begin{itemize}[nosep]
\item Shadow $+4.6\%$: cleaner prediction, no mass-calibration systematic
\item Internal theoretical consistency: field equation matches its sources
\item All local physics (P1--P13 patents): unaffected by $\psicosmo$
\end{itemize}
\end{tcolorbox}

%=============================================================================
\section{P14 Recommended Path Forward}
\label{sec:P14-forward}
%=============================================================================

\begin{enumerate}
\item \textbf{Adopt the EM-only interpretation as the baseline} for the
  cosmological sector. This resolves T1 and T2 and is self-consistent
  with the DFD field equation.

\item \textbf{Derive $a_0$ from the nonlinear self-coupling} or from the
  topological invariant, independently of $\psicosmo$. If $\kappa$ can be
  related to $\alpha^{57}$ or another DFD-intrinsic quantity, the MOND
  scale is recovered without Mach's principle.

\item \textbf{Develop a hybrid sourcing model}: the field equation
  (\ref{eq:field-eq}) currently couples $\psi$ to $T_{\rm EM}$.
  A self-consistent extension that includes a \emph{weak} matter coupling:
  \begin{equation}
  \Box\psi + (1+\kappa)(\nabla\psi)^2 = \frac{4\pi G}{c^4}
  \left[T_{\rm EM} + \epsilon\,T_{\rm matter}\right]
  \end{equation}
  with $\epsilon \ll 1$ could recover a small Machian contribution
  while keeping $\Gdot$ within bounds. The constraint from LLR
  requires $\epsilon \lesssim 3\ee{-3}$ (so that
  $\psicosmo^{\rm hybrid} = \epsilon\times 0.047 \lesssim 1.4\ee{-4}$).

\item \textbf{Pursue the shadow $+4.6\%$ as the leading P14 prediction}.
  This is independent of the cosmological sector resolution and should be
  the focus for next-generation EHT observations.
\end{enumerate}

%=============================================================================
\section{P15: Detection Roadmap with $M_{\rm eff}$ Correction}
\label{sec:P15}
%=============================================================================

\subsection{The W3i4 P11 correction: summary}

W3i4 P11 resolved the 28-order mode mass discrepancy. The placeholder
$\Mpsi = 10^{-6}$~kg used in W3i1--W3i3 was the mass of a hypothetical
oscillating membrane, not the effective inertial mass of the psi-mode.
The correct value from the DFD action principle is:
\begin{equation}
\Meff = \frac{\muG\,\Vcav}{4\pi G}
\label{eq:Meff-recall}
\end{equation}
giving $\Meff = 3.01\ee{6}$~kg for a TESLA cavity ($\Vcav = 3.84\ee{-3}$~m$^3$,
$\muG = 0.657$) and $\Meff^{\rm re-entrant} = 7.84\ee{5}$~kg for a
sub-litre re-entrant cavity.

\subsection{Impact on signal amplitude}

The psi-mode signal amplitude scales as $\Meff^{-1}$ (the driving force
per unit effective mass). The correction factor on signal amplitude:
\begin{equation}
\frac{q_{\rm corrected}}{q_{\rm W3i3}} = \frac{\Mpsi^{\rm used}}{\Meff}
= \frac{10^{-6}}{3.01\ee{6}} = 3.3\ee{-13}.
\label{eq:signal-correction}
\end{equation}

The SNR scales as signal/noise $\propto q/\sqrt{\text{noise}}$, so the
$|\lam-1|$ threshold (at SNR = 1) scales as $\sqrt{\Meff/\Mpsi^{\rm used}}$:
\begin{equation}
\frac{|\lam-1|_{\rm corrected}}{|\lam-1|_{\rm W3i3}}
= \sqrt{\frac{\Meff}{\Mpsi^{\rm used}}}
= \sqrt{\frac{3.01\ee{6}}{10^{-6}}} = \sqrt{3.01\ee{12}} = 1.74\ee{6}.
\label{eq:threshold-correction}
\end{equation}

All raw sensitivity thresholds derived from SQUID-detected psi-mode
amplitude using $\Mpsi = 10^{-6}$~kg must be multiplied by $1.74\ee{6}$.

\subsection{Which P15 reaches are affected?}

Following the W3i4 P11 analysis (Section~5), the impact depends on
\emph{how} each reach was derived.

\begin{table}[H]
\centering
\caption{P15 detection phases: which are affected by the $M_{\rm eff}$ correction.}
\label{tab:P15-affected}
\begin{tabular}{p{3cm}p{3cm}p{2.5cm}p{4cm}}
\toprule
\textbf{Phase} & \textbf{W3i4 Reach} & \textbf{$M_{\rm eff}$-Affected?}
  & \textbf{Derivation Basis} \\
\midrule
Phase 1 (CW) & $2\ee{-8}$ & \textbf{YES} & SQUID SNR from psi-mode amplitude \\
Phase 1+ (squeezed) & $3.6\ee{-9}$ & \textbf{YES} & Same, with squeezed noise floor \\
Phase 2 (pulsed) & $5.5\ee{-11}$ & \textbf{YES} & Pulsed energy / psi-mode coupling \\
Phase 2+ (pulsed+sq.) & $10^{-11}$ & \textbf{YES} & Same, with squeezing \\
Phase 3 (quantum net) & $3.7\ee{-12}$ & \textbf{PARTIAL} & Ratio from Phase~2, but anchored \\
Phase 3+ (q.+sq.) & $6.5\ee{-13}$ & \textbf{PARTIAL} & Same \\
P11+P2 joint & $8\ee{-13}$ & \textbf{NO} & Relative improvement ratio (W3i3) \\
ESS Lund & $1.14\ee{-14}$ & \textbf{NO} & $Q$-factor anomaly (passive) \\
LIGO 6th channel & $6\ee{-19}$ & \textbf{NO} & Optical path-length measurement \\
\bottomrule
\end{tabular}
\end{table}

\begin{finding}[Unaffected Reaches]
\label{find:unaffected}
Three P15/P11 detection reaches are \emph{not} affected by the $M_{\rm eff}$
correction:
\begin{enumerate}
\item \textbf{P11+P2 joint reach ($8\ee{-13}$)}: derived as a relative
  improvement ratio over the SQMS achieved bound, which is itself a measured
  quantity. $\Meff$ cancels in the ratio.
\item \textbf{ESS Lund ($1.14\ee{-14}$)}: derived from $Q$-factor anomaly
  detection in a passive mode, where the mode mass enters only through
  occupation number $n_\psi = 0$ (ground state). No $\Meff$ dependence.
\item \textbf{LIGO 6th channel ($6\ee{-19}$)}: optical path-length
  measurement; $\Meff$ does not appear.
\end{enumerate}
These three anchors are preserved. The P15 programme retains its ultimate
reach through these channels.
\end{finding}

\subsection{Corrected P15 Phase reaches}

For phases derived from raw SQUID SNR:

\begin{table}[H]
\centering
\caption{P15 detection phases: corrected for $M_{\rm eff} = 3.01\ee{6}$~kg.
The correction factor is $1.74\ee{6}$ on $|\lam-1|$ threshold.}
\label{tab:P15-corrected}
\begin{tabular}{lccc}
\toprule
\textbf{Phase} & \textbf{W3i4 Reach} & \textbf{Correction Factor}
  & \textbf{Corrected Reach} \\
\midrule
Phase 1 (CW, bare MZI) & $2\ee{-8}$ & $\times 1.74\ee{6}$ & $3.5\ee{-2}$ \\
Phase 1+ (squeezed) & $3.6\ee{-9}$ & $\times 1.74\ee{6}$ & $6.3\ee{-3}$ \\
Phase 2 (pulsed, $10^6$~J) & $5.5\ee{-11}$ & $\times 1.74\ee{6}$ & $9.6\ee{-5}$ \\
Phase 2+ (pulsed+squeezed) & $10^{-11}$ & $\times 1.74\ee{6}$ & $1.7\ee{-5}$ \\
\bottomrule
\end{tabular}
\end{table}

\begin{tcolorbox}[colback=red!5!white, colframe=red!60!black,
  title=\textbf{CRITICAL: Phases 1--2 Are Not Viable After $M_{\rm eff}$ Correction}]
The corrected Phase~1 CW reach is $|\lam-1| \sim 3.5\ee{-2}$. This
requires DFD coupling $\sim 10^7$ times larger than the current
experimental bound ($4.9\ee{-7}$). \textbf{Phase~1 cannot detect
any DFD signal within the allowed parameter space.}

Phase~1+ (squeezed): $6.3\ee{-3}$---still $\sim 10^4$ above the bound.

Phase~2 (pulsed, $10^6$~J): $9.6\ee{-5}$---still $\sim 200$ above the bound.

Phase~2+ (pulsed+squeezed): $1.7\ee{-5}$---$\sim 35$ above the bound.

\textbf{None of the Phases 1--2 can reach below the passive stability bound
$|\lam-1| < 4.9\ee{-7}$ after the $M_{\rm eff}$ correction.}
\end{tcolorbox}

\subsection{What survives: ratio-based and passive reaches}

\begin{table}[H]
\centering
\caption{P15/P11 detection reaches that \emph{survive} the $M_{\rm eff}$
correction, with status assessment.}
\label{tab:P15-surviving}
\begin{tabular}{llll}
\toprule
\textbf{Channel} & \textbf{Reach} & \textbf{TRL} & \textbf{Basis} \\
\midrule
P11+P2 joint & $8\ee{-13}$ & 2 & Relative ratio, SQMS-anchored \\
ESS Lund (passive $Q$) & $1.14\ee{-14}$ & 3 & $Q$-anomaly, no $\Meff$ \\
LIGO 6th channel & $6\ee{-19}$ & 2 & Optical path-length \\
Toroidal+YBCO (optimised) & $\sim 2\ee{-3}$ & 1 & SQUID SNR (corrected) \\
Levitated nanoparticle & See below & 1 & Different mass regime \\
\bottomrule
\end{tabular}
\end{table}

\subsection{Levitated nanoparticle detector: reassessment}

The W3i3 levitated nanoparticle detector (100~nm, 1~pg) was quoted as
achieving SNR~=~75/s at $|\lam-1| = 10^{-9}$. This detector has its own
effective mass $M_{\rm det} \sim 10^{-15}$~kg (the nanoparticle mass), not
the cavity $\Meff$. The psi-force on the nanoparticle scales as:
\begin{equation}
F_\psi = M_{\rm det}\,(\lam-1)\,\frac{c^2}{2}\,|\nabla\psi|.
\end{equation}
The signal here depends on $M_{\rm det}$, not $\Meff$, so the nanoparticle
detector reach is \emph{not} affected by the cavity mode mass correction
in the same way. However, the $\psi$-gradient at the detector location
is sourced by the cavity, so the gradient amplitude does depend on $\Meff$.

The corrected $\psi$-gradient from a cavity with $\Meff = 3.01\ee{6}$~kg
is reduced by $10^{-6}$ relative to W3i3 calculations. The nanoparticle
SNR is therefore also degraded by $\sim 10^{6}$, and the corrected reach
for the levitated nanoparticle is:
\begin{equation}
|\lam-1|_{\rm nano}^{\rm corrected} \sim 10^{-9}\times10^{6} = 10^{-3}.
\end{equation}
Still above the experimental bound. The nanoparticle ``canary'' detector
is not viable after the $\Meff$ correction.

\subsection{The material FOM table: still valid for relative ranking}

The W3i4 material FOM table ($\FOM_{\rm DFD} = \kappa\cdot\Tc^{1/2}$:
YBCO 965, MgB$_2$ 125, Nb$_3$Sn 92, Nb 3.2) is based on \emph{relative}
Process~A sourcing efficiency. Since $\Meff$ cancels in ratios between
materials at the same frequency and geometry, the FOM rankings are unaffected.

\begin{finding}[Material FOM Table Confirmed After $M_{\rm eff}$ Correction]
\label{find:FOM-valid}
The W3i4 practical material ranking (YBCO $>$ MgB$_2$ $>$ Nb$_3$Sn $>$ Nb)
is \emph{not affected} by the $M_{\rm eff}$ correction. The FOM quantifies
relative Process~A coupling efficiency, and $\Meff$ cancels in all ratios.
The Nb$_3$Sn+YBCO laminate recommendation stands.
\end{finding}

%=============================================================================
\section{Revised P15 Detection Strategy}
\label{sec:P15-strategy}
%=============================================================================

\subsection{The viable detection channels after $M_{\rm eff}$ correction}

With Phases 1--2 eliminated as viable direct detection routes, the P15
programme must pivot to the three surviving channels:

\begin{enumerate}
\item \textbf{ESS Lund passive $Q$-anomaly ($1.14\ee{-14}$)}: This is
  the tightest projected bound and is unaffected by $\Meff$. It relies on
  measuring anomalous $Q$-factor degradation in ESS SRF cavities during
  commissioning. The detection does not require exciting or measuring
  psi-mode amplitude; it detects the \emph{energy loss} to the psi-channel.
  \textbf{This becomes the primary near-term P15 strategy.}

\item \textbf{P11+P2 joint parametric detection ($8\ee{-13}$)}: Derived
  from a relative improvement ratio anchored to the SQMS achieved bound.
  Requires the purpose-built toroidal re-entrant cavity (P2) with the
  Nb$_3$Sn+YBCO laminate. TRL~2.

\item \textbf{LIGO 6th channel ($6\ee{-19}$)}: The most sensitive reach
  but requires LIGO-class interferometry with dedicated psi-detection
  modifications. TRL~2.
\end{enumerate}

\subsection{Revised P15 roadmap}

\begin{table}[H]
\centering
\caption{Revised P15 detection roadmap after $M_{\rm eff}$ correction.
Phases 1--2 are removed as non-viable.}
\label{tab:P15-revised}
\begin{tabular}{llcl}
\toprule
\textbf{Phase} & \textbf{Method} & $|\lam-1|_{\rm min}$ & \textbf{TRL} \\
\midrule
Phase A (near-term) & ESS Lund passive $Q$ & $1.14\ee{-14}$ & 3 \\
Phase B & P11+P2 toroidal re-entrant & $8\ee{-13}$ & 2 \\
Phase C & LIGO 6th channel & $6\ee{-19}$ & 2 \\
\bottomrule
\end{tabular}
\end{table}

\begin{tcolorbox}[colback=yellow!5!white, colframe=orange!60!black,
  title=\textbf{Revised P15 Strategy: Passive Detection Is the Only Near-Term Path}]
The $M_{\rm eff}$ correction eliminates all ``active'' SQUID-detected
psi-mode amplitude measurements as viable P15 routes. The surviving
channels are all \emph{passive} (detecting energy loss or optical
path-length changes, not psi-mode excitation amplitude).

\textbf{Key implication}: The SQMS hint at $|\lam-1| \sim 10^{-6}$
(W3i2) can no longer be probed by Phase~1 CW measurement. It can
only be tested by $Q$-anomaly measurements in high-$Q$ cavities
operated at or near the SQMS $Q$-frontier.

The ``discovery in 96 seconds'' headline from W3i3 (Phase~1 at SQMS
hint level) is \textbf{invalidated} by the $M_{\rm eff}$ correction.
\end{tcolorbox}

%=============================================================================
\section{Cross-References}
\label{sec:xrefs}
%=============================================================================

\subsection{P14 $\to$ P5 (MOND, Navigation, Timing)}

\begin{itemize}
\item P5 timing corrections (nonreciprocal delay, GPS corrections):
  \textbf{unaffected}. These depend on local $\nabla\psi$, not $\psicosmo$.
\item P5 MOND formula $a_0 = cH_0\psicosmo$: \textbf{invalid in EM-only
  framework}. Requires alternative MOND mechanism.
\item P5 Pioneer anomaly: \textbf{lost} unless alternative MOND mechanism
  preserves $a_0$.
\end{itemize}

\subsection{P14 $\to$ Cosmology module}

\begin{itemize}
\item Psi-tomography: \textbf{unaffected}. Does not require the Mach
  integral.
\item Dark energy as refractive bias: \textbf{lost} ($\psicosmo^{\rm EM}$
  too small by $10^5$).
\item S$_8$ prediction: \textbf{conditional} on alternative MOND mechanism.
\item Hubble tension: \textbf{conditional} on alternative MOND mechanism.
\item CMB Cold Spot (ISW): \textbf{conditional} on alternative MOND.
\item PTA spectral tilt: \textbf{conditional} on $\mu(x)$ survival.
\end{itemize}

\subsection{P15 $\to$ P2 (Resonators)}

\begin{itemize}
\item The toroidal re-entrant cavity (P2) becomes the critical enabling
  technology for P15 Phase~B. The $226\times$ improvement over TESLA
  baseline (W3i3) is a \emph{ratio} and survives the $M_{\rm eff}$
  correction.
\item The Nb$_3$Sn+YBCO laminate recommendation is confirmed (material
  FOM ratios unaffected).
\end{itemize}

\subsection{P15 $\to$ P11 (SRF Cavities)}

\begin{itemize}
\item The ESS Lund passive $Q$-anomaly detection (Phase~A) requires
  coordination with ESS commissioning schedule. This is now the
  \textbf{highest priority} P15/P11 experimental programme.
\item All P11 sensitivity thresholds from SQUID-detected amplitude are
  degraded by $1.74\ee{6}$ (same correction as P15 Phases 1--2).
\end{itemize}

%=============================================================================
\section{Summary of W3i5 Findings}
\label{sec:summary}
%=============================================================================

\begin{table}[H]
\centering
\caption{W3i5 P14+P15 Summary: all findings and their status.}
\label{tab:summary}
\small
\begin{tabular}{p{4.5cm}p{4cm}p{4cm}p{2cm}}
\toprule
\textbf{Item} & \textbf{W3i4 Status} & \textbf{W3i5 Status (EM-only)} & \textbf{Change} \\
\midrule
T1 ($\Gdot = 2.6\ee{-11}$) & Excluded $347\times$ & $7.3\ee{-17}$ (1000$\times$ below LLR) & \textbf{RESOLVED} \\
T2 (field eq.\ vs Mach) & Unresolved & Dissolved (EM-only self-consistent) & \textbf{RESOLVED} \\
P14+P5 incompatibility & $3.7\times$ mismatch & Dissolved by construction & \textbf{RESOLVED} \\
Mach's principle & $\psicosmo = 0.047$ & $\psicosmo^{\rm EM} = 1.3\ee{-7}$ (negligible) & \textbf{LOST} \\
MOND scale ($a_0$) & From $\psicosmo$ & Lost; self-coupling alternative proposed & \textbf{OPEN} \\
Dark energy (psi-screen) & Under DESI tension & Lost ($\psicosmo$ too small) & \textbf{LOST} \\
Shadow $+4.6\%$ & Compatible $0.3$--$0.8\sigma$ & Strengthened (no mass systematic) & \textbf{IMPROVED} \\
Hubble tension 70\% & Conditional on MOND & Conditional on alt.\ MOND & \textbf{CONDITIONAL} \\
S$_8$ prediction & Conditional on MOND & Conditional on alt.\ MOND & \textbf{CONDITIONAL} \\
\midrule
P15 Phase 1 (CW) & $2\ee{-8}$ & $3.5\ee{-2}$ (not viable) & \textbf{KILLED} \\
P15 Phase 1+ & $3.6\ee{-9}$ & $6.3\ee{-3}$ (not viable) & \textbf{KILLED} \\
P15 Phase 2 & $5.5\ee{-11}$ & $9.6\ee{-5}$ (not viable) & \textbf{KILLED} \\
P15 Phase 2+ & $10^{-11}$ & $1.7\ee{-5}$ (not viable) & \textbf{KILLED} \\
P15 Phase A (ESS passive) & $1.14\ee{-14}$ & $1.14\ee{-14}$ (unchanged) & \textbf{SURVIVES} \\
P11+P2 joint & $8\ee{-13}$ & $8\ee{-13}$ (unchanged) & \textbf{SURVIVES} \\
LIGO 6th channel & $6\ee{-19}$ & $6\ee{-19}$ (unchanged) & \textbf{SURVIVES} \\
Material FOM table & YBCO Rank 1 & YBCO Rank 1 (ratios unchanged) & \textbf{CONFIRMED} \\
\bottomrule
\end{tabular}
\end{table}

%=============================================================================
\section{Conclusion}
\label{sec:conclusion}
%=============================================================================

Wave~3 Iteration~5 demonstrates that the T1 $\Gdot$ falsification declared
in W3i4 can be \emph{completely resolved} by taking the DFD field equation
at face value: if $\psi$ is sourced only by electromagnetic stress-energy
(as Eq.~(\ref{eq:field-eq}) states), the cosmological psi-background from
all EM sources (CMB Thomson scattering, galactic and intergalactic magnetic
fields) is $\psicosmo^{\rm EM} \approx 1.3\ee{-7}$, yielding $\Gdot \approx
7.3\ee{-17}$~yr$^{-1}$---a factor of 1000 below the LLR bound. This
resolution eliminates T1, T2, and the P14+P5 incompatibility simultaneously.

The cost is significant: Mach's principle (strong form), the MOND scale
from $\psicosmo$, and the dark energy replacement via psi-screen refractive
bias are all lost. The MOND interpolating function $\mu(x)$ can potentially
be recovered through the nonlinear self-coupling fixed point, but this
requires $\kappa \sim 10^{26}$---an open theoretical question.

The shadow $+4.6\%$ prediction is \emph{strengthened} by the EM-only
interpretation (mass-calibration systematic eliminated) and remains the
cleanest testable DFD prediction in the cosmological sector.

For P15, the $M_{\rm eff} = 3.01\ee{6}$~kg correction eliminates all
SQUID-detected psi-mode amplitude measurements (Phases 1--2) as viable
detection routes. The surviving P15 channels are all passive or ratio-based:
ESS Lund $Q$-anomaly ($1.14\ee{-14}$), P11+P2 joint ($8\ee{-13}$), and
LIGO 6th channel ($6\ee{-19}$). The material FOM table (YBCO Rank~1) is
confirmed as ratios are unaffected.

\textbf{The central result of W3i5}: DFD has a self-consistent path that
resolves all three major cosmological tensions (T1, T2, P14+P5) at the
cost of the Machian programme and the cosmological-scale explanatory
power. The theory becomes \emph{narrower but internally consistent},
with the shadow $+4.6\%$ as the flagship cosmological prediction and
passive $Q$-anomaly detection as the primary experimental programme.

%=============================================================================
\end{document}
